\section{2026 Applied \& Computational Mathematics}

\begin{exercise}
\label{applied:2026:1}
Consider the equation
\[
\alpha \partial_t u(t,x)+\beta \partial_x u(t,x)-\gamma \partial_{xx}u(t,x)=f(x),
\]
for $(t,x)\in (0,T)\times (0,1)$, with $f\in L^2(0,1)$, $T>0$, boundary condition
$u(t,0)=u(t,1)=0$ for all $t>0$, and initial condition $u(0,x)=0$. The parameters satisfy
$\alpha>0$, $\beta\in \mathbb{R}$ and $\gamma>0$. Let $\mathcal{T}_h$ be a uniform mesh
partitioning $(0,1)$, namely the collection of intervals $[ih,(i+1)h]$ with
$i=0,1,\ldots,N$ and $h=1/(N+1)$.

(a) Write the fully discrete variational formulation of this problem using the continuous
piecewise linear finite element method (P1 Lagrange FEM) in space and the implicit Euler
method in time. Denote the time step by $\tau$ and $t_i=i\tau$ for all $i\in\mathbb{N}$.

(b) Prove the $L^2$ stability estimate $\|u_h^n\|\leq C\|f\|$ with a constant $C>0$
independent of $h$, $\tau$, and $n$.

(c) Let $\{\phi_i\}_{1\leq i\leq N}$ be the global Lagrange shape functions associated with
the nodes $x_i:=ih$ for $i=1,\ldots,N$. Denoting by
\[
u_h^i:=\sum_{1\leq j\leq N} U_j^i\phi_j
\]
the approximation of $u$ at $t_i$, write the algebraic linear system solved by
$(U_1^i,\ldots,U_N^i)$.
\end{exercise}

\begin{solution}
Let
\[
V_h=\{v_h\in C^0([0,1]): v_h|_K\in \mathbb{P}_1(K)\text{ for every }K\in\mathcal{T}_h,
\ v_h(0)=v_h(1)=0\}.
\]
The natural weak form is obtained by multiplying by $v_h\in V_h$ and integrating by parts
only in the diffusion term:
\[
\alpha(\partial_t u,v_h)+\beta(\partial_x u,v_h)+\gamma(\partial_x u,\partial_x v_h)
=(f,v_h),
\]
where $(p,q)=\int_0^1 p(x)q(x)\,dx$.

\textbf{(a) Fully discrete variational formulation.}
Set $u_h^0=0$. For each $n\geq 1$, find $u_h^n\in V_h$ such that, for every $v_h\in V_h$,
\[
\alpha\left(\frac{u_h^n-u_h^{n-1}}{\tau},v_h\right)
+\beta(\partial_x u_h^n,v_h)
+\gamma(\partial_x u_h^n,\partial_x v_h)
=(f,v_h).
\]
This is implicit Euler in time because all spatial operator terms are evaluated at the new
time level $t_n$.

\textbf{(b) $L^2$ stability.}
Take $v_h=u_h^n$. Since $u_h^n(0)=u_h^n(1)=0$,
\[
(\partial_x u_h^n,u_h^n)
=\frac{1}{2}\int_0^1 \partial_x\bigl((u_h^n)^2\bigr)\,dx
=0.
\]
Thus
\[
\frac{\alpha}{\tau}(u_h^n-u_h^{n-1},u_h^n)+\gamma\|\partial_x u_h^n\|^2=(f,u_h^n).
\]
Let $r_n=\|u_h^n\|$. The Cauchy--Schwarz inequality gives
\[
(u_h^n-u_h^{n-1},u_h^n)\geq r_n(r_n-r_{n-1}),
\]
and the Poincare inequality on $(0,1)$ gives $\|v\|\leq C_P\|\partial_x v\|$ for
$v\in H_0^1(0,1)$. Therefore, if $r_n>0$,
\[
\frac{\alpha}{\tau}(r_n-r_{n-1})+\frac{\gamma}{C_P^2}r_n\leq \|f\|.
\]
With $\mu=\gamma/C_P^2$, this implies
\[
r_n\leq \frac{\alpha}{\alpha+\mu\tau}r_{n-1}
+\frac{\tau}{\alpha+\mu\tau}\|f\|.
\]
Since $r_0=0$, induction yields
\[
\|u_h^n\|=r_n\leq \frac{1}{\mu}\|f\|
=\frac{C_P^2}{\gamma}\|f\|.
\]
Hence the required estimate holds with $C=C_P^2/\gamma$, independent of $h$, $\tau$, and
$n$.

\textbf{(c) Algebraic system.}
Write $U^n=(U_1^n,\ldots,U_N^n)^T$ and define
\[
M_{ij}=(\phi_j,\phi_i),\qquad
C_{ij}=(\phi_j',\phi_i),\qquad
K_{ij}=(\phi_j',\phi_i'),\qquad
F_i=(f,\phi_i).
\]
Testing the fully discrete equation by $\phi_i$ gives
\[
\left(\frac{\alpha}{\tau}M+\beta C+\gamma K\right)U^n
=\frac{\alpha}{\tau}M U^{n-1}+F.
\]
For the uniform P1 mesh, the standard entries are
\[
M_{ii}=\frac{2h}{3},\quad M_{i,i\pm 1}=\frac{h}{6},\qquad
K_{ii}=\frac{2}{h},\quad K_{i,i\pm 1}=-\frac{1}{h},
\]
with omitted entries equal to zero, and
\[
C_{ii}=0,\qquad C_{i,i+1}=\frac{1}{2},\qquad C_{i,i-1}=-\frac{1}{2}.
\]
\end{solution}

\begin{exercise}
\label{applied:2026:2}
Let $\Omega\subset \mathbb{R}^d$ be a bounded domain, $d\leq 3$. Consider the functional
$J:H_0^1(\Omega)\to \mathbb{R}$:
\[
J(u)=\int_\Omega \left(\frac{1}{2}|\nabla u(x)|^2+\frac{1}{4}u(x)^4-f(x)u(x)\right)\,dx,
\]
where $f\in L^2(\Omega)$.

(a) Compute the Frechet derivative (gradient) $\nabla J(u)$ and the second Frechet derivative
(Hessian) $\nabla^2J(u)$. Prove that $J$ is strictly convex.

(b) Write Newton's method for finding the minimizer $u^*$ of $J(u)$. Assuming the initial
guess $u_0$ is sufficiently close to the solution $u^*$, prove the second order convergence of
Newton's method:
\[
\|u_{k+1}-u^*\|_{H^1}\leq C\|u_k-u^*\|_{H^1}^2.
\]

(c) Let $s_k=u_{k+1}-u_k$ and $y_k=\nabla J(u_{k+1})-\nabla J(u_k)$. A Quasi-Newton method
resorts to an approximate Hessian $B_{k+1}$ satisfying the secant equation
$B_{k+1}s_k=y_k$. A BFGS update of $B_{k+1}$ is
\[
B_{k+1}v
=B_kv-\frac{\langle B_ks_k,v\rangle}{\langle B_ks_k,s_k\rangle}B_ks_k
+\frac{\langle y_k,v\rangle}{\langle y_k,s_k\rangle}y_k,
\]
for any test function $v\in H_0^1(\Omega)$. Assuming $B_k$ is positive definite, prove that the
above BFGS update is well-defined.
\end{exercise}

\begin{solution}
\textbf{(a) First and second variations.}
For $v,w\in H_0^1(\Omega)$,
\[
J'(u)[v]=\int_\Omega \nabla u\cdot \nabla v\,dx+\int_\Omega u^3v\,dx-\int_\Omega fv\,dx,
\]
and
\[
J''(u)[w,v]=\int_\Omega \nabla w\cdot \nabla v\,dx+3\int_\Omega u^2wv\,dx.
\]
The restriction $d\leq 3$ ensures the nonlinear terms are well-defined because
$H_0^1(\Omega)\hookrightarrow L^6(\Omega)$.

For $w\neq 0$,
\[
J''(u)[w,w]=\|\nabla w\|_{L^2}^2+3\int_\Omega u^2w^2\,dx>0.
\]
Thus $J$ has a positive definite second variation. Equivalently, the quadratic Dirichlet
energy is strictly convex on $H_0^1(\Omega)$, the $u^4$ term is convex, and the $-fu$ term is
linear. Hence $J$ is strictly convex and has at most one minimizer.

\textbf{(b) Newton iteration and local quadratic convergence.}
The minimizer $u^*$ satisfies
\[
J'(u^*)[v]=0\qquad \text{for all }v\in H_0^1(\Omega),
\]
or
\[
\int_\Omega \nabla u^*\cdot \nabla v\,dx+\int_\Omega (u^*)^3v\,dx
=\int_\Omega fv\,dx.
\]
Given $u_k$, Newton's method solves for the correction $\delta_k\in H_0^1(\Omega)$:
\[
J''(u_k)[\delta_k,v]=-J'(u_k)[v]\qquad \text{for all }v\in H_0^1(\Omega),
\]
and sets
\[
u_{k+1}=u_k+\delta_k.
\]
Equivalently, $u_{k+1}$ is characterized by
\[
\int_\Omega \nabla u_{k+1}\cdot \nabla v\,dx
+3\int_\Omega u_k^2u_{k+1}v\,dx
=\int_\Omega fv\,dx+2\int_\Omega u_k^3v\,dx.
\]

Let $F(u)=J'(u)$, interpreted as a map from $H_0^1(\Omega)$ to $H^{-1}(\Omega)$. Then
\[
F'(u)w=-\Delta w+3u^2w
\]
in weak form. The bilinear form
\[
a_u(w,v)=\int_\Omega \nabla w\cdot \nabla v\,dx+3\int_\Omega u^2wv\,dx
\]
is coercive because $a_u(w,w)\geq \|\nabla w\|_{L^2}^2$. By Lax--Milgram, $F'(u)$ is
invertible with a locally uniform bound on its inverse.

On bounded $H_0^1$ balls, $F'$ is Lipschitz as an operator $H_0^1\to H^{-1}$. Indeed,
using Sobolev embedding,
\[
\|(F'(u)-F'(v))w\|_{H^{-1}}
\leq C(\|u\|_{H^1}+\|v\|_{H^1})\|u-v\|_{H^1}\|w\|_{H^1}.
\]
For $e_k=u_k-u^*$, Newton's formula and $F(u^*)=0$ give
\[
e_{k+1}
=F'(u_k)^{-1}\bigl(F(u^*)-F(u_k)+F'(u_k)e_k\bigr).
\]
The Lipschitz estimate for $F'$ implies
\[
\|F(u^*)-F(u_k)+F'(u_k)e_k\|_{H^{-1}}
\leq C\|e_k\|_{H^1}^2.
\]
If $u_0$ is sufficiently close to $u^*$, all iterates remain in a neighborhood where
$\|F'(u_k)^{-1}\|$ is uniformly bounded, and therefore
\[
\|u_{k+1}-u^*\|_{H^1}\leq C\|u_k-u^*\|_{H^1}^2.
\]

\textbf{(c) Well-definedness of the BFGS update.}
Assume $s_k\neq 0$; if $s_k=0$, the secant update is unnecessary because the step has already
stopped. Since $B_k$ is positive definite,
\[
\langle B_ks_k,s_k\rangle>0.
\]
It remains to check $\langle y_k,s_k\rangle>0$. By the fundamental theorem of calculus in
Banach spaces,
\[
\langle y_k,s_k\rangle
=\int_0^1 J''(u_k+ts_k)[s_k,s_k]\,dt.
\]
Using the expression for the Hessian,
\[
\langle y_k,s_k\rangle
=\int_0^1\left(\|\nabla s_k\|_{L^2}^2
+3\int_\Omega (u_k+ts_k)^2s_k^2\,dx\right)\,dt>0.
\]
Both denominators in the BFGS formula are therefore strictly positive. Hence the update is
well-defined for every nonzero step $s_k$.
\end{solution}

\begin{exercise}
\label{applied:2026:3}
Consider the initial value problem over $\mathbb{R}^N$ in the form
\[
x'=f(t,x),\qquad x(0)=x_0\in \mathbb{R}^N,
\]
where $f:[0,T]\times \mathbb{R}^N\to \mathbb{R}^N$ is smooth. Consider the family of one-step
methods
\[
x_{n+1}=x_n+(1-b)hf(t_n,x_n)+bhf(t_{n+1},x_{n+1}),
\]
where $h=t_{n+1}-t_n$ for any $n\geq 0$ is a uniform step size and $b\in[0,1]$ is a constant.

(a) Find the value of $b$ so that the local truncation error is $O(h^3)$.

(b) Apply the method to $x'=\lambda x$, $x(0)=x_0\in \mathbb{R}$. Find the function $g(\cdot)$
such that $x_n=g(h\lambda)^n x_0$.

(c) Determine the values of $b$ so that this method is A-stable.
\end{exercise}

\begin{solution}
\textbf{(a) Local truncation error.}
Let $x(t)$ be the exact solution and suppress the time $t_n$ in the notation. The one-step
defect obtained by inserting the exact solution is
\[
x(t+h)-x(t)-h\bigl((1-b)x'(t)+bx'(t+h)\bigr).
\]
Taylor expansion gives
\[
x(t+h)=x(t)+hx'(t)+\frac{h^2}{2}x''(t)+\frac{h^3}{6}x'''(t)+O(h^4),
\]
and
\[
x'(t+h)=x'(t)+hx''(t)+\frac{h^2}{2}x'''(t)+O(h^3).
\]
Hence the defect equals
\[
\left(\frac{1}{2}-b\right)h^2x''(t)
+\left(\frac{1}{6}-\frac{b}{2}\right)h^3x'''(t)+O(h^4).
\]
For the local truncation error to be $O(h^3)$, the $h^2$ term must vanish, so
\[
b=\frac{1}{2}.
\]

\textbf{(b) Stability function.}
For $x'=\lambda x$, set $z=h\lambda$. The method becomes
\[
x_{n+1}=x_n+(1-b)zx_n+bzx_{n+1}.
\]
Thus
\[
(1-bz)x_{n+1}=(1+(1-b)z)x_n,
\]
and
\[
x_n=g(z)^n x_0,\qquad
g(z)=\frac{1+(1-b)z}{1-bz}.
\]

\textbf{(c) A-stability.}
The method is A-stable exactly when
\[
|g(z)|\leq 1\qquad \text{for all }\operatorname{Re}z\leq 0.
\]
Writing $z=x+iy$, the inequality $|1+(1-b)z|\leq |1-bz|$ is equivalent to
\[
|1-bz|^2-|1+(1-b)z|^2=-2x+(2b-1)|z|^2\geq 0.
\]
If $b\geq 1/2$, then $-2x\geq 0$ for $\operatorname{Re}z=x\leq 0$ and
$(2b-1)|z|^2\geq 0$, so the whole left half-plane is contained in the stability region.

If $b<1/2$, take $z=iy$ with $|y|$ large. Then
\[
-2\operatorname{Re}z+(2b-1)|z|^2=(2b-1)y^2<0,
\]
so the imaginary axis is not contained in the stability region. Therefore the method is
A-stable precisely for
\[
b\in\left[\frac{1}{2},1\right].
\]
\end{solution}

\begin{exercise}
\label{applied:2026:4}
Let $A$ be an invertible $N\times N$ matrix. The shifted QR iteration for a given sequence of
shifts $\{\sigma_n\}$ is defined by
\[
A_0=A,\qquad A_n-\sigma_n I_N=Q_nR_n,\qquad
A_{n+1}=R_nQ_n+\sigma_n I_N,
\]
where $I_N$ is the $N\times N$ identity matrix, $Q_n$ is orthogonal, and $R_n$ is upper
triangular with positive diagonal entries.

(a) Prove that if no $\sigma_n$ is an eigenvalue of $A$, then the sequences $\{A_n\}$,
$\{Q_n\}$, and $\{R_n\}$ are uniquely defined and satisfy
\[
A_{n+1}=Q_n^T A_n Q_n,\qquad A_{n+1}=R_nA_nR_n^{-1}.
\]

(b) Suppose $A$ is a symmetric $2\times 2$ matrix with eigenvalues $\lambda_1$ and
$\lambda_2$. Let $\sigma_0=\lambda_1$. Find $A_1$.

(c) Let $\widehat{Q}_n=Q_0\cdots Q_{n-1}$ and
$\widehat{R}_n=R_{n-1}\cdots R_0$ for $n\geq 1$. Prove
\[
\widehat{Q}_{k+1}\widehat{R}_{k+1}
=\prod_{i=0}^k (A-\sigma_i I_N).
\]
\end{exercise}

\begin{solution}
\textbf{(a) Similarity identities and uniqueness.}
Assume inductively that $A_n$ is orthogonally similar to $A$. Then $A_n$ and $A$ have the
same eigenvalues. Since $\sigma_n$ is not an eigenvalue of $A$, the matrix
$A_n-\sigma_n I_N$ is nonsingular. Every nonsingular square matrix has a unique QR
factorization with $Q_n$ orthogonal and $R_n$ upper triangular with positive diagonal entries.
Thus $Q_n$ and $R_n$ are uniquely defined, and so is $A_{n+1}$.

From $A_n-\sigma_n I_N=Q_nR_n$, we have $A_n=Q_nR_n+\sigma_n I_N$. Therefore
\[
Q_n^TA_nQ_n
=Q_n^T(Q_nR_n+\sigma_n I_N)Q_n
=R_nQ_n+\sigma_n I_N
=A_{n+1}.
\]
This proves that $A_{n+1}$ is orthogonally similar to $A_n$, hence to $A$, and closes the
induction. Also,
\[
R_nA_nR_n^{-1}
=R_n(Q_nR_n+\sigma_n I_N)R_n^{-1}
=R_nQ_n+\sigma_n I_N
=A_{n+1}.
\]

\textbf{(b) Exact eigenvalue shift.}
Here $\sigma_0=\lambda_1$ is an eigenvalue of $A$, so the hypothesis in part (a) fails:
$A-\lambda_1I$ is singular and the QR factorization with all diagonal entries of $R$ positive
does not exist. Thus, under the literal definition stated in the problem, $A_1$ is not
well-defined.

If $\lambda_1\neq \lambda_2$ and one uses the standard rank-deficient extension of the QR
step, the result is the familiar one-step deflation. For an unreduced symmetric $2\times 2$
matrix, $A-\lambda_1I$ has rank one. The first QR vector can be chosen in the range of
$A-\lambda_1I$, which is the eigenspace for $\lambda_2$, while the second QR vector spans the
nullspace of $A-\lambda_1I$, the eigenspace for $\lambda_1$. With this convention,
\[
A_1=Q_0^TAQ_0=
\begin{pmatrix}
\lambda_2 & 0\\
0 & \lambda_1
\end{pmatrix}.
\]
The important point is that this conclusion uses an extra convention for the singular QR
step; it is not forced by the positive-diagonal nonsingular QR definition in part (a).

\textbf{(c) Product formula.}
Assume the QR steps under discussion are defined, for example under the hypothesis of part
(a).
First note from part (a) that
\[
A_n=\widehat{Q}_n^T A\widehat{Q}_n.
\]
We prove the product identity by induction. For $k=0$,
\[
\widehat{Q}_1\widehat{R}_1=Q_0R_0=A-\sigma_0I_N.
\]
Assume
\[
\widehat{Q}_k\widehat{R}_k=\prod_{i=0}^{k-1}(A-\sigma_iI_N).
\]
Since
\[
Q_kR_k=A_k-\sigma_kI_N
=\widehat{Q}_k^T(A-\sigma_kI_N)\widehat{Q}_k,
\]
we obtain
\begin{align*}
\widehat{Q}_{k+1}\widehat{R}_{k+1}
&=\widehat{Q}_k Q_kR_k\widehat{R}_k\\
&=\widehat{Q}_k\widehat{Q}_k^T(A-\sigma_kI_N)\widehat{Q}_k\widehat{R}_k\\
&=(A-\sigma_kI_N)\prod_{i=0}^{k-1}(A-\sigma_iI_N).
\end{align*}
All factors are polynomials in $A$, hence commute. Therefore
\[
\widehat{Q}_{k+1}\widehat{R}_{k+1}
=\prod_{i=0}^k(A-\sigma_iI_N).
\]
\end{solution}

\begin{exercise}
\label{applied:2026:5}
Let $A\in \mathbb{R}^{n\times n}$ be an $n$ by $n$ real matrix and let $\sigma_i(A)$ be its
$i$-th largest singular value. A vector $x\in \mathbb{R}^n$ such that $Ax=x$ is called a fixed
point of $A$.

(a) Assume $\sigma_1(A)\leq 1$. Show that every fixed point of $A$ is a fixed point of its
transpose $A^T$.

(b) Consider
\[
A=
\begin{pmatrix}
1 & 1\\
0 & 0
\end{pmatrix}
\]
to verify that the assertion in (a) is not generally true without assuming $\sigma_1(A)\leq 1$.

(c) Assume $AA^T=A^TA$. Show that every fixed point of $A$ is a fixed point of its transpose
$A^T$.
\end{exercise}

\begin{solution}
\textbf{Statement-risk note.}
The statement is consistent as written: part (b) is explicitly asking for a
counterexample when the singular-value hypothesis in (a) is removed, while (c)
uses normality as a different sufficient condition.

\textbf{(a) Contraction case.}
The assumption $\sigma_1(A)\leq 1$ means $\|A\|_2\leq 1$, or equivalently
$I-A^TA$ is positive semidefinite. Let $Ax=x$. Then
\[
\|x\|_2^2=\|Ax\|_2^2=x^TA^TAx\leq x^Tx=\|x\|_2^2.
\]
Equality holds, so
\[
x^T(I-A^TA)x=0.
\]
Because $I-A^TA$ is positive semidefinite, this implies
\[
(I-A^TA)x=0.
\]
Using $Ax=x$,
\[
A^Tx=A^TAx=x.
\]
Thus every fixed point of $A$ is also a fixed point of $A^T$.

\textbf{(b) Counterexample without the singular-value assumption.}
For
\[
A=
\begin{pmatrix}
1 & 1\\
0 & 0
\end{pmatrix},
\]
the fixed point equation $Ax=x$ gives
\[
\begin{pmatrix}
x_1+x_2\\
0
\end{pmatrix}
=
\begin{pmatrix}
x_1\\
x_2
\end{pmatrix},
\]
so $x_2=0$. Hence every vector $(a,0)^T$ is a fixed point of $A$. But
\[
A^T
\begin{pmatrix}
a\\
0
\end{pmatrix}
=
\begin{pmatrix}
1 & 0\\
1 & 0
\end{pmatrix}
\begin{pmatrix}
a\\
0
\end{pmatrix}
=
\begin{pmatrix}
a\\
a
\end{pmatrix},
\]
which equals $(a,0)^T$ only when $a=0$. Thus a nonzero fixed point of $A$ need not be a fixed
point of $A^T$. Indeed, $\sigma_1(A)=\sqrt{2}>1$.

\textbf{(c) Normal real matrix case.}
Let $B=A-I$. The condition $AA^T=A^TA$ implies
\[
BB^T=(A-I)(A^T-I)=AA^T-A-A^T+I
\]
and
\[
B^TB=(A^T-I)(A-I)=A^TA-A^T-A+I.
\]
Thus $BB^T=B^TB$, so $B$ is normal. For a real normal matrix $B$,
\[
\|Bx\|_2^2=x^TB^TBx=x^TBB^Tx=\|B^Tx\|_2^2.
\]
If $Ax=x$, then $Bx=(A-I)x=0$. The identity above gives $\|B^Tx\|_2=0$, hence
\[
B^Tx=(A^T-I)x=0.
\]
Therefore $A^Tx=x$.
\end{solution}

\begin{exercise}
\label{applied:2026:6}
Let the subdifferential of a convex function $f:\mathbb{R}^n\to \mathbb{R}\cup\{+\infty\}$
at a point $x$ be denoted by $\partial f(x)$. Recall that a vector $x^*\in \mathbb{R}^n$ is a
subgradient of $f$ at $x$ if it satisfies
\[
f(z)\geq f(x)+\langle x^*,z-x\rangle,\qquad \forall z\in \mathbb{R}^n.
\]
Namely,
\[
\partial f(x)=\{x^*: f(z)\geq f(x)+\langle x^*,z-x\rangle,\ \forall z\in \mathbb{R}^n\}.
\]
Here, $\langle x,y\rangle=x^Ty$ for any vectors $x$ and $y$ in $\mathbb{R}^n$. We assume
$n\geq 2$; the $n=1$ case is trivial.

(a) Let $f:\mathbb{R}^n\to \mathbb{R}$ be defined by
\[
f(x)=\max_{i=1,\ldots,n}x_i,\qquad x=(x_1,\ldots,x_n)^T\in \mathbb{R}^n.
\]
Compute the subdifferential $\partial f(0)$.

(b) Let
\[
g(x)=\max_{i=1,\ldots,n}x_i+\delta_{\mathbb{R}_+^n}(x),
\]
where $\delta_{\mathbb{R}_+^n}$ is the indicator function of the non-negative orthant
$\mathbb{R}_+^n$, namely $\delta_{\mathbb{R}_+^n}(x)=0$ if $x\geq 0$ and $+\infty$ otherwise.
Show that
\[
\partial g(0)=\partial f(0)-\mathbb{R}_+^n.
\]
Here, one may use the fact that
$\partial g(x)=\partial f(x)+\partial \delta_{\mathbb{R}_+^n}(x)$ for
$g(x)=f(x)+\delta_{\mathbb{R}_+^n}(x)$.
\end{exercise}

\begin{solution}
\textbf{(a) Subdifferential of the maximum at the origin.}
The function can be written as
\[
f(x)=\max_{1\leq i\leq n}\langle e_i,x\rangle,
\]
where $e_i$ is the $i$-th standard basis vector. At $x=0$, every affine function
$\langle e_i,x\rangle$ is active. The subdifferential of a pointwise maximum of finitely many
linear functions is the convex hull of the active gradients, so
\[
\partial f(0)=\operatorname{conv}\{e_1,\ldots,e_n\}.
\]
Equivalently,
\[
\partial f(0)=\left\{p\in \mathbb{R}^n:\ p_i\geq 0\text{ for all }i,
\sum_{i=1}^n p_i=1\right\}.
\]
To check this directly, if $p$ is in the simplex, then
\[
\langle p,z\rangle=\sum_{i=1}^n p_i z_i\leq \max_i z_i=f(z),
\]
and since $f(0)=0$, this is exactly the subgradient inequality at $0$. Conversely, the
subgradient inequality with $z=t\mathbf{1}$ and
$t>0$ gives $\sum_i p_i\leq 1$, while $t<0$ gives $\sum_i p_i\geq 1$, hence
$\sum_i p_i=1$. Taking $z=-te_j$ with $t>0$ gives $0=f(z)\geq -tp_j$, so $p_j\geq 0$ for
each $j$.

\textbf{(b) Adding the non-negative orthant constraint.}
The subdifferential of the indicator function is the normal cone. At the origin,
\[
\partial \delta_{\mathbb{R}_+^n}(0)
=N_{\mathbb{R}_+^n}(0)
=\{q\in \mathbb{R}^n:\langle q,z\rangle\leq 0\text{ for all }z\in \mathbb{R}_+^n\}.
\]
This cone is exactly the non-positive orthant:
\[
N_{\mathbb{R}_+^n}(0)=-\mathbb{R}_+^n.
\]
Using the stated sum rule,
\[
\partial g(0)
=\partial f(0)+\partial \delta_{\mathbb{R}_+^n}(0)
=\partial f(0)-\mathbb{R}_+^n.
\]
\end{solution}
